% Owner: member A (table) and member D. Target length: 1.25 pages.
% Source: docs/feature_handling.md.
\section{Preprocessing and Feature Engineering}\label{sec:preprocessing}

We removed 425 exact duplicate training rows, which leaves 149{,}575 training rows. No test rows were removed. Town, block and street names were normalised before the join with the HDB block table, which raised the share of matched rows from about 79\% to 100\%.

Table~\ref{tab:feature-handling} lists the decision for every column of the original data.

\begin{table}[htbp]
\caption{Handling of each column in the original data}
\label{tab:feature-handling}
\centering
\footnotesize
\begin{tabular}{p{0.24\columnwidth}p{0.66\columnwidth}}
\hline
\textbf{Column} & \textbf{Decision and reason} \\
\hline
Approval date & Transformed into a month index (months since January 2021) and a cyclic month. Defines the time split. \\
Town & Kept as a categorical feature (26 levels). \\
Block, street & Used only to build the address key for the block join. Not model inputs, because of high cardinality (2{,}720 and 1{,}146 levels), and because coordinates and subzone carry the location after the join. \\
Flat type & Spellings unified. Kept as a categorical and as an ordinal (1-room to executive). \\
Flat model & Kept as a categorical feature (20 levels). \\
Floor area & Kept, plus its logarithm and the area per room. \\
Lease year & Kept, plus the flat age at approval and a flag for one invalid row. \\
Furnished, fee & Dropped. Constant in train and test. \\
Monthly rent & Target. Not transformed in the stored data. \\
\hline
\end{tabular}
\end{table}

\todo{Derived location features (distances and counts for MRT, schools, malls, rail lines, campuses) and the external data (SingStat CPI, COE, URA rail geometry). Two short paragraphs.}

\todo{The about 50 suspicious rows (large flats with very low rent): what they are and why we kept them. Member C's ablation shows that removing them does not help.}
